\clearpage
\fancyhead[L]{\small BEOBACHTUNGEN.}
\SourceLeaf{469}
\begin{center}
\rule{0.55\linewidth}{0.4pt}\\[2pt]
Astronomische Nachrichten. Band XVI. Nr. 378. Seite 303. 1839 Mai 30.\\[2pt]
\rule{0.55\linewidth}{0.4pt}
\end{center}

Herr Hofrath \textsc{Gauss} hat in Göttingen nur den Anfang dieser Sonnenfinsterniss um
\begin{center}1838 März 15.\quad $3^{\rm u}59^{\rm m}9^{\rm s}6$ mittlere Zeit\end{center}
\noindent beobachten können. Von der Mitte der Finsterniss an wurde die Sonne durch Wolken unsichtbar gemacht.
\begin{flushright}\textsc{Schumacher.}\end{flushright}

\clearpage
\hypertarget{native-469-occultation}{}
\begin{center}
\rule{0.55\linewidth}{0.4pt}\\[2pt]
Astronomische Nachrichten. Band XVIII. Nr. 430. Seite 367. 1841 August 5.\\[2pt]
\rule{0.55\linewidth}{0.4pt}
\end{center}

Herr Hofrath \textsc{Gauss} hat am 23. Mai 1841 in Göttingen, da das ungünstige Wetter die Beobachtung des Eintritts verhinderte, den Austritt von $42\omega'$ Gemin. um $9^{\rm u}32^{\rm m}10^{\rm s}6$ m. Z. beobachtet. Herr Dr. \textsc{Goldschmidt} beobachtete diesen Austritt $0^{\rm s}3$ früher.
\begin{flushright}\textsc{Schumacher.}\end{flushright}

\clearpage
\hypertarget{native-469-notice1842}{}
\begin{center}
\rule{0.55\linewidth}{0.4pt}\\[2pt]
Astronomische Nachrichten. Band XX. Nr. 465. Seite 160. 1842 Januar 5.\\[2pt]
\rule{0.55\linewidth}{0.4pt}
\end{center}

[Mittheilung der von \textsc{Goldschmidt} ausgeführten Beobachtungen des \emph{Cometen.}]

\bigskip
\hypertarget{native-469-notice1844}{}
\begin{center}
\rule{0.55\linewidth}{0.4pt}\\[2pt]
Astronomische Nachrichten. Band XXI. Nr. 494. Seite 221. 222. 1844 Januar 18.\\[2pt]
\rule{0.55\linewidth}{0.4pt}
\end{center}

[Mittheilung der von Dr. \textsc{Goldschmidt} berechneten elliptischen Elemente des Cometen.]

\clearpage
\hypertarget{native-469-letter}{}
\begin{center}
\rule{0.55\linewidth}{0.4pt}\\[2pt]
Astronomisehe Nachrichten. Band XXI. Nr. 495. S. 235..238. 1844 Februar 8.\\[2pt]
\rule{0.55\linewidth}{0.4pt}
\end{center}

\begin{center}[\textsc{Gauss} an \textsc{Schumacher.}]\quad Göttingen 1844 Januar 17.\end{center}

Anbei übersende ich Ihnen eine Ephemeride des \emph{Cometen}, welche Hr. Dr. \textsc{Goldschmidt} nach seinen zweiten Elementen berechnet hat. Ich zweifle nicht, dass man ihn damit immer leicht wird auffinden können, wenn nicht das gar zu schwach werdende Licht des Cometen die Sichtbarkeit hindert.

\Needspace{3\baselineskip}
\SourceLeaf{470}
Eine Vergleichung dieser zweiten Elemente mit den sämmtlichen ihm bekannt gewordenen Beobachtungen wird Hr. Dr. \textsc{Goldschmidt} unverzüglich beendigen.

Meine eigene Beobachtung vom 11. Januar steht so:
\begin{center}
\begin{tabular}{l|l|l|l}
1844 Januar 11 & $9^{\rm u}23^{\rm m}41^{\rm s}.8$ & $77^\circ8'49''4$ & ---\quad ---\quad --- \\
 & $9\;23\;47.5$ & ---\quad ---\quad --- & $+3^\circ32'26''7$ \\
\end{tabular}
\end{center}

Die Declination bleibt etwas ungewiss. Der geringe Unterschied zwischen den Zeiten, für welche die Gerade Aufsteigung und die Abweichung gelten, rührt daher, dass jene sich auf drei, letztere auf zwei Beobachtungen gründet, und dass die dritte blos für Gerade Aufsteigung brauchbare Beobachtung \emph{nahe}, aber \emph{nicht genau}, mitten zwischen den beiden andern liegt.

Ich habe jetzt die Chronometer von einem Hannoverschen Künstler hier zur Prüfung, jedoch erst seit 4 oder 5 Tagen. In dieser Zeit ist der Gang vortrefflich gleichförmig, obwohl eine so kurze Zeit noch nicht viel entscheiden kann.

Schliesslich bemerke ich noch, dass meine Beobachtung vom 13. December in der neuen Reduction eine kleine Aenderung erlitten hat, aber nur in Folge eines bei der \emph{Sternreduction} vorgefallenen Irrthums. Sie steht jetzt so:
\begin{center}
\begin{tabular}{l|l|l|l}
1842 December 13 & $9^{\rm u}35^{\rm m}15^{\rm s}.9$ & ---\quad ---\quad --- & $+3^\circ39'28''6$ \\
 & $9\;43\;13.7$ & $78^\circ45'38''0$ & ---\quad ---\quad --- \\
\end{tabular}
\end{center}

\clearpage
\hypertarget{native-470-report}{}
\begin{center}
\rule{0.55\linewidth}{0.4pt}\\[2pt]
Astronomische Nachrichten. Band XXI. Nr. 496. Seite 247. 1844 Februar 17.\\[2pt]
\rule{0.55\linewidth}{0.4pt}
\end{center}

Herr Hofrath \textsc{Gauss} hat mir in einem Briefe vom 16. folgende Beobachtungen des Cometen gesandt:
\begin{center}
\begin{tabular}{l|l|l|l}
1844 Januar 15 & $8^{\rm u}18^{\rm m}17^{\rm s}8$ m. Z. & AR.: $77^\circ20'59''6$ & $\delta$: ---\quad ---\quad --- \\
 & $8\;10\;44.5$\quad --- & AR.: ---\quad ---\quad --- & $\delta$: $+3^\circ49'47''3$ \\
\end{tabular}
\end{center}

Der Comet ging einem Sterne, dessen Position für diesen Tag Dr. \textsc{Goldschmidt} zu $78^\circ17'2''3$ $+3^\circ51'2''9$ berechnet hatte um $8^{\rm u}18^{\rm m}17^{\rm s}8$ vor $3'34''18$, und war um $8^{\rm u}10^{\rm m}44^{\rm s}5$ m. Z. $1'15''6$ südlicher. Herr Hofrath \textsc{Gauss} hält diese Beobachtung für so gut, als er sie bei der grossen Lichtschwäche des Cometen mit dem Kreismicrometer machen konnte, wenn, wie es sich von selbst versteht, der Stern gut bestimmt ist.

In einem spätern Brief vom 21. hat Herr Hofrath \textsc{Gauss} mir gefälligst die Vergleichung mit den zweiten \textsc{Goldschmidt}schen Elementen mitgetheilt.
\begin{flushright}\textsc{Schumacher.}\end{flushright}
